%% titre: Quiz — Espaces vectoriels
%% niveau: Licence 1
%% matiere: Algèbre linéaire
%% annee: 2025-2026
%% description: Exemple de QCM (à remplacer), corrigé automatiquement.

\begin{question}
Soit $E=\{(x,y)\in\R^2 \mid x+y=0\}$. L'ensemble $E$ est :
\begin{reponses}
  \bonne un sous-espace vectoriel de $\R^2$
  \mauvaise une partie de $\R^2$ qui n'est pas stable par addition
  \mauvaise l'ensemble vide
\end{reponses}
\begin{explication}
$E=\Ker\big((x,y)\mapsto x+y\big)$ est le noyau d'une application linéaire.
\end{explication}
\end{question}

\begin{question}
Quelles familles sont libres dans $\R^2$ ? (plusieurs réponses possibles)
\begin{reponses}
  \bonne $\{(1,0),(0,1)\}$
  \mauvaise $\{(1,2),(2,4)\}$
  \bonne $\{(1,1),(1,-1)\}$
  \mauvaise $\{(1,0),(0,1),(1,1)\}$
\end{reponses}
\begin{explication}
Dans $\R^2$, une famille libre a au plus $2$ vecteurs, et $(2,4)=2\,(1,2)$ donne une famille liée.
\end{explication}
\end{question}

\begin{question}
Si $A\in\mathcal{M}_n(\R)$ vérifie $A^2=A$, alors $\rg A + \rg(I_n-A)$ vaut :
\begin{reponses}
  \mauvaise $n-1$
  \bonne $n$
  \mauvaise $2n$
\end{reponses}
\end{question}
